\chapter{Conclusions}

This dissertation has developed an operator-valued theory of dynamic grid
strength in converter-rich power systems and used it to define dynamic hosting
capacity, virtual-inertia placement, weak elements, damping-aware network
reinforcement, and a traceable source-to-claim ledger for the project
artifacts. The hierarchy begins with the operational graph induced by the
power-flow operating point, continues through commuting fixed damping, isolates
the uniform-decay case \(D=\gamma M\), passes through fixed non-commuting
damping, continues to rational self-energy \(\Pi(s)\) after controller-state
condensation, introduces the Green function \(T(s)^{-1}\) as the carrier of
finite-action behavior, and then separates genuine nonlinear memory from
small-signal NEP theory. The central conclusion is that a damping, hosting, or
grid-strength claim is meaningful only relative to the operator, port/action
family, margin, and validation gate that carry the mechanism.

\section{Theoretical Conclusions}

A first conclusion is that graph topology becomes dynamic only after the
operating point is included. In the lossless retained model, the linearized
stiffness is
\[
  L(z^\star)=B W(z^\star)B^T,\qquad
  w_e=\kappa_e\cos\eta_e .
\]
Thus line loading, angle cohesion, cuts, cycles, inertia placement, and local
control fields enter the same graph-time operator. Static graph metrics are
useful screens, but the protected object is the response of
\(T(s;z^\star)=s^2M+sD+BWB^T+\Pi(s;z^\star)\).

The second conclusion is that \(D=\gamma M\) is a special floor, not the full
decouplable class. The exact fixed-\(D\) decoupling condition is
\([\Lt,\Dt]=0\). Under \(D=\gamma M\), graph modes decouple and every
underdamped mode has real part \(-\gamma/2\). Under Rayleigh or other
commuting damping, graph modes still decouple, but modal decay is not uniform.
This distinction gives a precise foundation for classical graph-frequency
intuition and also limits it.

The third conclusion is that fixed non-commuting damping breaks modal
decoupling. The graph eigenvectors diagonalize \(\Lt\), but the modal damping
matrix \(\Gamma=Q^T\Dt Q\) need not be diagonal. Off-diagonal damping creates
coupled modal dynamics, and eliminating the other modes produces rational
corrections. Thus, even a fixed-\(D\) system can require a frequency-dependent
modal interpretation when one insists on reducing to a single retained mode.

The fourth conclusion is that controller-state condensation changes the reduced
operator class. A Schur complement of the eliminated controller block gives
\[
  T(s)=s^2M+sD_0+L+\Pi(s),
  \qquad
  \Pi(s)=\Pi_K(s)+s\Pi_D(s).
\]
This rational correction has poles and residues inherited from the eliminated
dynamics. A constant damping matrix cannot represent those features exactly.
Static damping bridges are therefore local approximations whose validity must
be tested, not assumed.

The fifth conclusion is that dynamic effective resistance closes the loop
between graph metrics and finite actions. The Green function
\(G_T(s)=T(s)^{-1}\) defines
\[
  \mathcal R_e(s)=a_e^TG_T(s)a_e ,
\]
which reduces to classical effective resistance at the static graph limit and
becomes a meromorphic port compliance away from zero. A finite line action
satisfies \(1+\Delta b_e\mathcal R_e(s)=0\), and general low-rank actions
satisfy a small action determinant. Therefore local sensitivities are tangent
approximations of exact finite-action laws, not standalone planning objects.

The sixth conclusion is that reduced models need certificates. A QEP screen,
static self-energy approximation, data-identified \(\widehat\Pi(s)\), or
condensed NEP should be accepted over a protected contour only when the
relative perturbation bound
\[
  \sup_{\Gamma}\|\widehat T(s)^{-1}(T(s)-\widehat T(s))\|_2<1
\]
or an equivalent certificate preserves the protected eigenvalue count. A
margin derivative should be treated as a recommendation only when its sign is
larger than its error bound.

The seventh conclusion is that fragility is structured. An unstructured
pseudospectrum warns about non-normality, but a planning decision needs a
physical perturbation family: line uncertainty, damping, inertia, controller
gains, delays, operating-point variation, or identified-model error. The
resulting fragility radius measures the smallest admissible perturbation that
can move a pole to a protected boundary.

The eighth conclusion is that NEP theory and nonlinear stability are different
modules. The Schur NEP preserves linearized memory around an operating point.
The nonlinear DAE layer must still address equilibrium regularity,
bifurcations, finite-amplitude stability, and regions of attraction. A
small-signal damping margin is therefore a local gate, not a large-signal
certificate.

The ninth conclusion is that hosting capacity in converter-rich systems is
not only a voltage or thermal number. A static-feasible IBR penetration can
still violate a damping-ratio margin, a protected resolvent-gain threshold, or
a controller-resonance gate. The first-order spectral hosting screen is
therefore a useful tangent screen, but it becomes an inner certificate only with
curvature bounds and must be checked against the QEP or NEP when modal coupling
or controller self-energy is important.

The tenth conclusion is that graph reinforcement and dynamic robustness are
not synonymous. A line upgrade may improve algebraic connectivity or effective
resistance while reducing damping headroom. The correct planning object is a
vector containing graph gain, damping-margin derivative, protected-gain
derivative, feasibility, and possible Braess-like or control-resonant reversal
flags.

\section{Consequences for Research Practice}

The hierarchy imposes discipline on future simulation and planning studies.
Before a graph metric is used as evidence, the study must state whether the
operator is an operational graph at a fixed operating point, uniform
proportional damping, commuting fixed damping, non-commuting fixed damping,
rational self-energy, or nonlinear memory. Before a weak bus or weak line is
named, the margin and action family must be named. Before a reduced converter
model is trusted, the spectral effect of removed states must be visible in
either the retained matrix or the self-energy.

This thesis therefore separates three epistemic levels:

\begin{enumerate}
\item a certificate, which implies a margin under explicit assumptions;
\item a screen, which identifies candidates but does not imply the margin;
\item a gate, which tests whether the screen is reliable enough for the next
claim.
\end{enumerate}

The distinction is not editorial. It is necessary because the same graph
quantity can be a certificate under proportional damping, a screen under fixed
non-commuting damping, and an incomplete descriptor after controller
condensation.

\section{Open Theoretical Work}

Several theoretical problems remain open.

\begin{enumerate}
\item Develop graph-time resonance atlases that relate operating-point
geometry, graph-mode conversion, controller self-energy, and protected
frequency-response peaks.
\item Turn the total derivative of \(L=B W(z^\star)B^T\) into a continuation
tool for planning actions, so that line, dispatch, and converter changes are
distinguished from frozen-stiffness perturbations.
\item Develop computable error bounds between the full Schur NEP and static or
identified approximations of \(\Pi(s)\) or its damping-channel projection
\(\Sigma(s)\) over specified spectral regions.
\item Turn the dynamic effective resistance law
\(1+\Delta b_e\mathcal R_e(s)=0\) into scalable algorithms for finite line
screening and certified reinforcement budgets.
\item Define non-commutation measures tailored to inertia-normalized power
networks, with explicit links to damping-ratio error.
\item Connect passivity or positive-real properties of \(\Sigma(s)\) to
small-signal damping margins in non-normal retained systems.
\item Replace scalar passivity summaries with directional passivity matrices
or dynamic multipliers when MIMO channel geometry is important.
\item Develop a passivity-retuning protocol in which converter-controller
gains are adjusted to improve a band-limited self-energy passivity index while
also moving the protected Schur NEP pole in a damping-improving direction.
\item Compute structured fragility radii for topology, line loading, damping,
inertia, controller, delay, and operating-point uncertainty families.
\item Characterize when grid-strength or Laplacian placement metrics are
certificates, screens, or invalid descriptors for a given converter-control
family.
\item Test whether cellular-sheaf Laplacians with bus/controller stalks and
line restriction maps reveal weak nodes or weak links that scalar Laplacian
sensitivity, fixed-\(D\) QEP screens, and Schur self-energy diagnostics miss.
\item Build theory-to-validation protocols that combine NEP sensitivity,
pole-residue analysis, and full DAE eigenvalue recomputation.
\item Combine Loewner or contour-integral NEP extraction with continuation
tracking so that local pole-motion maps can be followed through high-dimensional
hosting and controller-retuning sweeps.
\item Extend the dynamic hosting screen beyond first order, including
non-commuting damping, controller self-energy, nonlinear memory, structured
fragility, and uncertainty in operating points.
\item Identify \(\Pi(s)\) from black-box converter data with stability,
causality, port, and passivity constraints, then propagate identification
error to certified planning decisions.
\item Register line-reinforcement and virtual-inertia planning baselines so
that Braess-like reversals and placement gains can be compared against
meaningful controls.
\end{enumerate}

\section{Closing Statement}

The transition from synchronous-machine dominated grids to converter-rich
grids is often described as a loss of inertia. That description is incomplete.
The deeper theoretical change is a change in the strength, damping, fragility,
and memory operator. Uniform viscous damping, commuting fixed damping,
non-commuting modal damping, rational controller self-energy, Green-function
action calculus, structured fragility, and nonlinear memory are different
mathematical objects. Static hosting, dynamic hosting, virtual inertia,
controller retuning, and line reinforcement are likewise different planning
questions unless they share the same port, operator, and protected margin. A
doctoral theory of dynamic grid strength must keep those objects separate and
then reconnect them through certificates. This dissertation provides that path:
from static topology to operational graph-time resonance, from effective
resistance to dynamic effective resistance, from fixed \(D\) to \(\Pi(s)\),
from \(\Pi(s)\) to structured fragility, and from linear self-energy to
nonlinear memory.
